Fix integers $d\ge 1$ and $n\ge 2$. For a finite set $P\subset \mathbb{R}^d$ of $m$ \emph{distinct} points, write
\[D(P)=\{\|x-y\| : x,y\in P,\ x\neq y\}\]
for the set of (Euclidean) distances determined by $P$.
We say that $Q\subseteq P$ with $|Q|=n$ is \emph{pairwise-distance-distinct} if the $\binom{n}{2}$ distances
\[\{\|x-y\|: x,y\in Q,\ x\neq y\}
\]
are all distinct (equivalently, $|D(Q)|=\binom{n}{2}$).
Define $f_d(n)$ to be the least integer $m$ such that every $m$-point set $P\subset \mathbb{R}^d$ contains a pairwise-distance-distinct subset $Q$ of size $n$.

Edge cases/conventions:
\begin{itemize}
\item $f_d(2)=2$.
\item For $n\ge 3$, the condition means every triangle (when $n=3$) is scalene, and more generally every two distinct edges in the complete graph on $Q$ have different lengths.
\end{itemize}
